\documentclass[10pt,a4paper]{article}

% --- Кодировка, шрифты и русский язык ---
\usepackage{cmap}
\usepackage[T2A]{fontenc}
\usepackage[utf8]{inputenc}
\usepackage[english,russian]{babel}
\usepackage{microtype}

% --- Математика ---
\usepackage{amsmath,amssymb,amsfonts,mathtools}

% --- Геометрия и оформление ---
\usepackage[a4paper,left=15mm,right=15mm,top=18mm,bottom=18mm,headheight=14pt]{geometry}
\usepackage[dvipsnames,svgnames]{xcolor}
\usepackage{enumitem}
\usepackage{booktabs}
\usepackage{array}
\usepackage{tabularx}
\usepackage{fancyhdr}
\usepackage{titlesec}
\usepackage[most]{tcolorbox}
\usepackage[colorlinks=true,linkcolor=NavyBlue,urlcolor=NavyBlue,citecolor=NavyBlue]{hyperref}

% --- Цвета и блоки ---
\definecolor{AccentBlue}{HTML}{1B365D}
\definecolor{SubtleBg}{HTML}{F4F6F9}
\definecolor{BorderBlue}{HTML}{B8C7DC}

\newtcolorbox{summarybox}[1][]{
  enhanced,
  breakable,
  colback=SubtleBg,
  colframe=BorderBlue,
  boxrule=0.6pt,
  arc=2pt,
  left=6pt,right=6pt,top=5pt,bottom=5pt,
  fonttitle=\bfseries\color{AccentBlue},
  coltitle=AccentBlue,
  colbacktitle=SubtleBg,
  attach title to upper={.\ },
  #1
}

% --- Настройка заголовков и списков ---
\titleformat{\section}
  {\Large\bfseries\color{AccentBlue}}
  {\thesection.}{0.5em}{}[\vspace{2pt}\titlerule]
\titlespacing*{\section}{0pt}{12pt}{6pt}

\titleformat{\subsection}
  {\large\bfseries\color{AccentBlue!90!black}}
  {\thesubsection.}{0.5em}{}
\titlespacing*{\subsection}{0pt}{9pt}{4pt}

\titleformat{\subsubsection}
  {\normalsize\bfseries}
  {\thesubsubsection.}{0.5em}{}
\titlespacing*{\subsubsection}{0pt}{6pt}{2pt}

\setlist[itemize]{leftmargin=1.5em,itemsep=2pt,topsep=3pt,parsep=0pt,beginpenalty=10000}
\setlist[enumerate]{leftmargin=1.7em,itemsep=2pt,topsep=3pt,parsep=0pt,beginpenalty=10000}
\setlength{\parindent}{0pt}
\setlength{\parskip}{3pt}

% --- Колонтитулы ---
\pagestyle{fancy}
\fancyhf{}
\fancyhead[L]{\footnotesize\textcolor{gray}{TeamBook: Тервер и комбинаторика}}
\fancyhead[R]{\footnotesize\textcolor{gray}{\nouppercase{\leftmark}}}
\fancyfoot[C]{\small\thepage}
\renewcommand{\headrulewidth}{0.4pt}

% --- Математические макросы и русская типографика ---
\renewcommand{\le}{\leqslant}
\renewcommand{\ge}{\geqslant}
\renewcommand{\emptyset}{\varnothing}

\newcommand{\N}{\mathbb{N}}
\newcommand{\Z}{\mathbb{Z}}
\newcommand{\R}{\mathbb{R}}
\newcommand{\F}{\mathbb{F}}
\newcommand{\E}{\mathbb{E}}
\newcommand{\Var}{\operatorname{Var}}
\newcommand{\Ind}{\mathbb{I}}
\newcommand{\iverson}[1]{[\,#1\,]}
\newcommand{\fall}[2]{#1^{\underline{#2}}}
\newcommand{\rise}[2]{#1^{\overline{#2}}}
\newcommand{\stirlingfirst}[2]{\genfrac{[}{]}{0pt}{}{#1}{#2}}
\newcommand{\stirlingsecond}[2]{\genfrac{\{}{\}}{0pt}{}{#1}{#2}}
\DeclareMathOperator{\lcm}{lcm}
\DeclareMathOperator{\rk}{rank}
\DeclareMathOperator{\inv}{inv}
\DeclareMathOperator{\Fix}{Fix}
\DeclareMathOperator{\Stab}{Stab}
\DeclareMathOperator{\sgn}{sgn}
\DeclareMathOperator{\popcnt}{popcount}
\DeclareMathOperator{\Geom}{Geom}
\DeclareMathOperator{\Unif}{Unif}

\begin{document}
\thispagestyle{plain}

% --- Заголовок ---
\begin{center}
  {\LARGE\bfseries\color{AccentBlue} Теория вероятностей и комбинаторика}\\[4pt]
  {\normalsize Краткий практический справочник по теории для олимпиадной информатики}
\end{center}
\vspace{2pt}

\begin{summarybox}[title={Базовые обозначения и соглашения}]
\begin{itemize}
  \item $\N = \{1, 2, \dots\}$ — натуральные числа, $\N_0 = \{0, 1, 2, \dots\}$ — целые неотрицательные, $\Z$ — все целые, $\R$ — вещественные; $[n] = \{1, 2, \dots, n\}$ — первые $n$ натуральных чисел ($[0] = \emptyset$).
  \item $|A|$ — размер (число элементов) конечного множества $A$; $A \sqcup B$ — объединение непересекающихся множеств ($A \cap B = \emptyset$); $A \triangle B = (A \setminus B) \cup (B \setminus A)$ — элементы, лежащие ровно в одном из $A, B$.
  \item $\iverson{P}$ — \textbf{скобка Айверсона}: равна $1$, если условие $P$ выполнено, и $0$, если нет. По соглашению $0! = 1$ и $0^0 = 1$.
  \item $d \mid n$ — число $n$ нацело делится на $d$ ($d \nmid n$ — не делится); $\gcd(a, b)$ и $\lcm(a, b)$ — НОД и НОК.
  \item $\fall{x}{k} = x(x-1)\dots(x-k+1)$ — \textbf{убывающий факториал} (произведение $k$ чисел вниз от $x$); $\rise{x}{k} = x(x+1)\dots(x+k-1)$ — \textbf{возрастающий факториал} ($\fall{x}{0} = \rise{x}{0} = 1$).
  \item $S_n$ — множество всех $n!$ перестановок чисел $\{1, \dots, n\}$; запись $[x^n] F(x)$ означает «взять коэффициент при $x^n$ в многочлене или степенном ряде $F(x)$».
\end{itemize}
\end{summarybox}

{\small\tableofcontents}
\vspace{4pt}
\hrule
\vspace{6pt}

% ==============================================================================
\section{Уровень 1: База и модульная арифметика}
% ==============================================================================

\subsection{Модульная арифметика и базовая комбинаторика}

\textbf{Сравнения и остатки по модулю.} Запись $a \equiv b \pmod m$ ($m \ge 2$) означает, что целые числа $a$ и $b$ дают одинаковый остаток при делении на $m$ (то есть разность $a - b$ делится на $m$), а $(a \bmod m) \in \{0, 1, \dots, m - 1\}$ — сам этот остаток. Складывать, вычитать и умножать числа можно сразу по модулю $m$ на каждом шаге, не дожидаясь конца вычислений. Делить по модулю напрямую нельзя — деление заменяется умножением на обратный элемент.

\textbf{Бинарное возведение в степень (\texttt{binpow}).} Позволяет вычислить $a^n \bmod m$ ($n \in \N_0$) за $O(\log n)$ умножений, на каждом шаге уменьшая степень вдвое при чётном $n$ и отщепляя один множитель при нечётном:
\[
  a^0 = 1, \qquad
  a^n \equiv
  \begin{cases}
    (a^{n/2})^2 \pmod m, & n \equiv 0 \pmod 2, \\
    a \cdot a^{n-1} \pmod m, & n \equiv 1 \pmod 2.
  \end{cases}
\]

\textbf{Обратный элемент по модулю.} Обратным к числу $a$ по модулю $m$ называется такой остаток $a^{-1} \in \{1, \dots, m - 1\}$, что $a \cdot a^{-1} \equiv 1 \pmod m$ (умножение на $a^{-1}$ работает как «деление на $a$»). Обратный элемент существует тогда и только тогда, когда $\gcd(a, m) = 1$:
\begin{itemize}
  \item \textbf{Простой модуль $p$ ($p \nmid a$):} по малой теореме Ферма $a^{p-1} \equiv 1 \pmod p$, поэтому $a^{-1} \equiv a^{p-2} \pmod p$ — находится одним \texttt{binpow} за $O(\log p)$.
  \item \textbf{Произвольный модуль $m$ ($\gcd(a, m) = 1$):} либо через расширенный алгоритм Евклида за $O(\log m)$ (находим целые $x, y$ из уравнения $ax + my = 1$, тогда $a^{-1} = x \bmod m$), либо по теореме Эйлера $a^{-1} \equiv a^{\varphi(m)-1} \pmod m$, где $\varphi(m)$ — количество чисел от $1$ до $m$, взаимно простых с $m$.
  \item \textbf{Линейный предподсчёт всех обратных $1^{-1}, \dots, N^{-1} \bmod p$} ($p$ простое, $N < p$) за $O(N)$ без логарифма:
  \[
    \inv[1] = 1, \qquad \inv[i] = \bigl(p - \lfloor p / i \rfloor \cdot \inv[p \bmod i] \bmod p\bigr) \bmod p \quad (i = 2, \dots, N).
  \]
\end{itemize}

\textbf{Линейный предподсчёт \texttt{fact[]} и \texttt{inv\_fact[]}.} Чтобы отвечать на запросы биномиальных коэффициентов за $O(1)$ по простому модулю $p$ (при $N < p$), предподсчитывают массив факториалов $\mathrm{fact}[i] = i! \bmod p$ слева направо, берут один обратный элемент от $N!$ за $O(\log p)$, а все остальные $\mathrm{inv\_fact}[i] = (i!)^{-1} \bmod p$ досчитывают справа налево (так как $\frac{1}{(i-1)!} = \frac{i}{i!}$) суммарно за $O(N + \log p)$:
\[
  \mathrm{fact}[0] = 1,\quad \mathrm{fact}[i] = \mathrm{fact}[i-1] \cdot i \bmod p \;\;(1 \le i \le N);
\]
\[
  \mathrm{inv\_fact}[N] = (\mathrm{fact}[N])^{p-2} \bmod p,\quad
  \mathrm{inv\_fact}[i-1] = \mathrm{inv\_fact}[i] \cdot i \bmod p \;\;(i = N, N-1, \dots, 1).
\]
Отсюда же за $O(1)$ выражается одиночный обратный элемент: $i^{-1} \equiv \mathrm{fact}[i-1] \cdot \mathrm{inv\_fact}[i] \pmod p$.

\textbf{Базовые комбинаторные объекты} (выбираем элементы из множества размера $n$):
\begin{itemize}
  \item \textbf{Правила суммы и произведения:} если нужно выбрать объект \emph{либо} из $A$, \emph{либо} из $B$ ($A \cap B = \emptyset$), число способов равно $|A| + |B|$; если нужно последовательно выбрать первый элемент из $A$ \emph{и} независимо второй из $B$ — число пар равно $|A| \cdot |B|$.
  \item \textbf{Перестановки ($P_n$):} число способов выстроить все $n$ различных элементов в ряд (порядок важен): $P_n = n! = 1 \cdot 2 \dots n$.
  \item \textbf{Перестановки с повторениями (мультиномиальный коэффициент):} число способов выстроить в ряд $n$ предметов, среди которых $k_1$ одинаковых предметов $1$-го типа, $\dots$, $k_m$ одинаковых $m$-го типа ($k_1 + \dots + k_m = n$), либо разбить $n$ различных элементов по $m$ пронумерованным коробкам фиксированных размеров $k_1, \dots, k_m$:
  \[
    \binom{n}{k_1, k_2, \dots, k_m} = \frac{n!}{k_1!\,k_2!\dots k_m!}.
  \]
  \item \textbf{Размещения без повторений ($A_n^k$) и с повторениями ($\bar{A}_n^k$):} число упорядоченных наборов длины $k$ из $n$ элементов (порядок выбора важен). Без повторений (все $k$ элементов различны) и с повторениями (элементы можно брать многократно — это просто слова длины $k$ над алфавитом размера $n$):
  \[
    A_n^k = \fall{n}{k} = \frac{n!}{(n - k)!} \quad (0 \le k \le n, \text{ иначе } 0), \qquad \bar{A}_n^k = n^k.
  \]
  \item \textbf{Сочетания без повторений ($C_n^k = \binom{n}{k}$):} число способов выбрать подмножество из $k$ элементов среди $n$ различных, когда порядок внутри выборки \emph{не важен}:
  \[
    \binom{n}{k} = \frac{n!}{k!\,(n - k)!} = \frac{\fall{n}{k}}{k!} \quad (0 \le k \le n), \qquad \binom{n}{k} = 0 \text{ при } k < 0 \text{ или } k > n.
  \]
  По простому модулю $p > n$: $\binom{n}{k} \equiv \mathrm{fact}[n] \cdot \mathrm{inv\_fact}[k] \cdot \mathrm{inv\_fact}[n-k] \pmod p$ за $O(1)$.
\end{itemize}

\textbf{Бином Ньютона и рабочие тождества.} Раскрытие скобок $(a + b)^n$ для любого $n \in \N_0$:
\[
  (a + b)^n = \sum_{k=0}^{n} \binom{n}{k} a^{n-k} b^k.
\]
Самые частые свойства для упрощения сумм:
\begin{itemize}
  \item \textbf{Симметрия и суммы по строке:} $\binom{n}{k} = \binom{n}{n-k}$; $\sum_{k=0}^n \binom{n}{k} = 2^n$ (все подмножества); $\sum_{k=0}^n (-1)^k \binom{n}{k} = \iverson{n = 0}$ (чётных и нечётных подмножеств поровну при $n \ge 1$); $\sum_{k=0}^n k\binom{n}{k} = n 2^{n-1}$.
  \item \textbf{Внесение множителя (поглощение):} $k\binom{n}{k} = n\binom{n-1}{k-1}$ ($k \ge 1$); $\binom{n}{m}\binom{m}{k} = \binom{n}{k}\binom{n-k}{m-k}$.
  \item \textbf{Свёртка Вандермонда} (выбор $k$ элементов из двух куч размеров $n$ и $m$): $\sum_{j=0}^{k} \binom{n}{j}\binom{m}{k-j} = \binom{n+m}{k}$; частный случай: $\sum_{j=0}^n \binom{n}{j}^2 = \binom{2n}{n}$.
\end{itemize}

\textbf{Треугольник Паскаля.} Рекуррентная таблица значений $\binom{n}{k}$ (каждый элемент либо берём в подмножество, либо не берём). Главный плюс — считает $\binom{n}{k} \bmod M$ за $O(N^2)$ по \textbf{абсолютно любому} (в том числе составному) модулю $M$, так как использует только сложение:
\[
  \binom{n}{0} = \binom{n}{n} = 1, \qquad \binom{n}{k} = \binom{n-1}{k-1} + \binom{n-1}{k} \quad (1 \le k \le n - 1).
\]

\textbf{Тождество хоккейной клюшки (Hockey-Stick Identity).} Сумма биномиальных коэффициентов вдоль диагонали треугольника Паскаля сворачивается в одну клетку «на загибе клюшки» ($n, r, m \in \N_0$):
\[
  \sum_{i=r}^{n} \binom{i}{r} = \binom{n+1}{r+1} \quad (\text{при } n < r \text{ сумма равна } 0), \qquad
  \sum_{j=0}^{m} \binom{r+j}{j} = \binom{r+m+1}{m}.
\]

% ------------------------------------------------------------------------------
\subsection{Шары и перегородки (Stars and Bars)}

\textbf{Идея метода.} Разложить $n$ одинаковых (неразличимых) шаров по $k \ge 1$ различным ящикам — то же самое, что выстроить в ряд $n$ шаров и $k - 1$ перегородку между ящиками, либо найти число целых решений уравнения $x_1 + x_2 + \dots + x_k = n$ ($x_i$ — сколько шаров попало в $i$-й ящик):
\begin{itemize}
  \item \textbf{Пустые ящики разрешены ($x_i \ge 0$, сочетания с повторениями $\bar{C}_k^n$):} выбираем $k - 1$ место под перегородки из $n + k - 1$ символов:
  \[
    |\{(x_1, \dots, x_k) \in \Z^k : x_1 + \dots + x_k = n,\; x_i \ge 0\}| = \binom{n + k - 1}{k - 1} = \binom{n + k - 1}{n}.
  \]
  \item \textbf{В каждом ящике хотя бы один шар ($x_i \ge 1$):} ставим $k - 1$ перегородку в $n - 1$ промежуток между шарами (чтобы перегородки не слипались):
  \[
    |\{(x_1, \dots, x_k) \in \Z^k : x_1 + \dots + x_k = n,\; x_i \ge 1\}| = \binom{n - 1}{k - 1} \quad (n \ge k, \text{ иначе } 0).
  \]
\end{itemize}

\textbf{Сдвиг при произвольных ограничениях снизу ($x_i \ge a_i$).} Если в $i$-й ящик нужно заранее положить минимум $a_i \in \Z$ шаров, делаем замену $y_i = x_i - a_i \ge 0$. Получаем базовую задачу $y_1 + \dots + y_k = n'$, где $n' = n - \sum_{i=1}^k a_i$ — остаток свободных шаров. Ответ: $\binom{n' + k - 1}{k - 1}$ при $n' \ge 0$ (и $0$ при $n' < 0$).

\textbf{Переход от неравенств к равенствам.} Чтобы решить неравенство $x_1 + \dots + x_k \le n$ ($x_i \ge 0$), заводим $(k + 1)$-й «мусорный ящик» $x_{k+1} = n - (x_1 + \dots + x_k) \ge 0$, куда скидываем все неиспользованные шары. Получаем уравнение $x_1 + \dots + x_k + x_{k+1} = n$ ($x_i \ge 0$), число его решений равно $\binom{n + k}{k}$. Строгое неравенство $< n$ для целых чисел заменяется на $\le n - 1$.

\textbf{Ограничения сверху через ФВИ ($a_i \le x_i \le b_i$).} Сначала сдвигом приводим нижние границы к нулю: $x_1 + \dots + x_k = n$, $0 \le x_i \le c_i$ ($c_i = b_i - a_i \ge 0$). Назовём $i$-й ящик «переполненным» (плохое свойство $A_i$), если в нём оказалось $x_i \ge c_i + 1$ шаров. Перебирая подмножество $S \subseteq [k]$ заведомо переполненных ящиков и заранее откладывая в каждый $i \in S$ по $c_i + 1$ шаров, по формуле включений-исключений получаем:
\[
  \sum_{S \subseteq [k]} (-1)^{|S|} \binom{n - \sum_{i \in S}(c_i + 1) + k - 1}{k - 1}
  \quad \bigl(\text{слагаемое равно } 0, \text{ если верхний индекс } < k - 1\bigr).
\]
Если верхние ограничения одинаковы ($0 \le x_i \le c$), все множества одного размера $j = |S|$ дают равный вклад, и сумма считается за $O\bigl(\min(k, n / (c + 1))\bigr)$:
\[
  \sum_{j=0}^{\min(k,\, \lfloor n / (c + 1) \rfloor)} (-1)^j \binom{k}{j} \binom{n - j(c + 1) + k - 1}{k - 1}.
\]

% ------------------------------------------------------------------------------
\subsection{Разложение Кантора и код Лемера}

\textbf{Код Лемера.} Для перестановки $\pi = (\pi_1, \dots, \pi_n) \in S_n$ её \textbf{кодом Лемера} называется массив $L(\pi) = (c_1, c_2, \dots, c_n)$, где $c_i$ показывает, сколько чисел \emph{правее} позиции $i$ меньше $\pi_i$ (то есть сколько ещё не использованных элементов меньше того, который мы поставили на $i$-е место):
\[
  c_i = |\{j > i : \pi_j < \pi_i\}|, \qquad 0 \le c_i \le n - i \quad (c_n = 0).
\]
Код Лемера однозначно задаёт перестановку (это биекция), а сумма $\sum_{i=1}^n c_i = \inv(\pi)$ равна общему числу инверсий в $\pi$.

\textbf{Факториальная система счисления (разложение Кантора).} Любое целое число $X \in [0, n! - 1]$ единственным способом записывается по «базису факториалов», где цифра на позиции $i$ не превосходит $n - i$:
\[
  X = \sum_{i=1}^{n} c_i \cdot (n - i)!, \qquad 0 \le c_i \le n - i \quad (c_n = 0).
\]
Чтобы перевести число $X$ в цифры $(c_1, \dots, c_n)$ с конца, последовательно для $k = 1, 2, \dots, n$ берём $c_{n + 1 - k} = X \bmod k$ и делим $X \leftarrow \lfloor X / k \rfloor$.

\textbf{Номер перестановки по порядку.} Если подставить код Лемера перестановки $\pi$ в разложение Кантора, число $\rk(\pi) = \sum_{i=1}^n c_i (n - i)! \in [0, n! - 1]$ в точности равно лексикографическому номеру перестановки $\pi$ (в $0$-нумерации: тождественная перестановка $(1, 2, \dots, n)$ имеет номер $0$, а $(n, \dots, 2, 1)$ — номер $n! - 1$).

\textbf{Быстрый перевод туда и обратно за $O(n \log n)$ через дерево Фенвика.} Храним $1$-индексированное дерево Фенвика $T[1 \dots n]$ из нулей и единиц: $1$ означает, что число ещё свободно, $0$ — уже поставлено в перестановку (в начале все $n$ элементов равны $1$):
\begin{itemize}
  \item \textbf{Перестановка $\to$ код Лемера:} идём слева направо ($i = 1, \dots, n$). Число свободных элементов меньше $\pi_i$ — это просто сумма на префиксе $[1, \pi_i - 1]$ в Фенвике: $c_i = \mathrm{sum}(\pi_i - 1)$. Записываем $c_i$ и помечаем число использованным: $\mathrm{add}(\pi_i, -1)$.
  \item \textbf{Код Лемера $\to$ перестановка:} идём слева направо ($i = 1, \dots, n$). Нам нужно выбрать $(c_i + 1)$-е по счёту свободное число — это минимальный индекс $v$, где префиксная сумма в $T$ достигает $c_i + 1$. Чтобы не тратить $O(\log^2 n)$ на бинпоиск, делаем спуск по степеням двойки прямо внутри массива дерева Фенвика за $O(\log n)$: ставим $\mathrm{pos} = 0$, $\mathrm{rem} = c_i + 1$; для $b = 2^{\lfloor \log_2 n \rfloor}, \dots, 2, 1$, если $\mathrm{pos} + b \le n$ и $T[\mathrm{pos} + b] < \mathrm{rem}$, шагаем $\mathrm{pos} \mathrel{+}= b$ и вычитаем $\mathrm{rem} \mathrel{-}= T[\mathrm{pos}]$. Тогда $\pi_i = \mathrm{pos} + 1$, после чего делаем $\mathrm{add}(\pi_i, -1)$. Общее время — $O(n \log n)$.
\end{itemize}

% ------------------------------------------------------------------------------
\subsection{Базовый тервер}

\textbf{Исходы, события и вероятность — базовый словарь.}
\begin{itemize}
  \item \textbf{Случайный эксперимент и элементарный исход ($\omega$):} любой процесс со случайным результатом (бросок костей, генерация случайного графа). Один полностью конкретизированный неделимый вариант того, чем закончился опыт, называется \emph{элементарным исходом} $\omega$, а множество вообще всех возможных исходов обозначается $\Omega$ (\emph{пространство исходов}).
  \item \textbf{Событие ($A \subseteq \Omega$):} любое высказывание об исходе опыта, задаваемое просто набором подходящих исходов $A \subseteq \Omega$ (говорят, что событие $A$ \emph{наступило}, если выпавший исход $\omega$ принадлежит $A$).
  \begin{itemize}
    \item $\Omega$ — \emph{достоверное событие} (наступает всегда), $\emptyset$ — \emph{невозможное событие} (не наступает никогда), $\bar{A} = \Omega \setminus A$ — \emph{противоположное событие} («$A$ не произошло»).
    \item $A \cap B$ — произошли \textbf{одновременно оба} события $A$ и $B$; $A \cup B$ — произошло \textbf{хотя бы одно} из $A, B$.
    \item События называют \textbf{несовместными}, если они не могут произойти вместе ($A \cap B = \emptyset$).
  \end{itemize}
  \item \textbf{Вероятность события ($P(A) \in [0, 1]$):} в дискретных задачах у каждого элементарного исхода $\omega \in \Omega$ есть свой неотрицательный вес (шанс выпадения) $P(\{\omega\})$, причём сумма весов по всему $\Omega$ равна $1$. Вероятность события $A$ — сумма весов всех входящих в него исходов: $P(A) = \sum_{\omega \in A} P(\{\omega\})$. Базовые правила:
  \[
    P(\Omega) = 1, \qquad P(\emptyset) = 0, \qquad P(\bar{A}) = 1 - P(A), \qquad
    P(A \cup B) = P(A) + P(B) - P(A \cap B)
  \]
  (для попарно несовместных событий пересечений нет и вероятности просто складываются: $P(\bigsqcup_i A_i) = \sum_i P(A_i)$).
\end{itemize}

\textbf{Классическая вероятность.} Если число возможных исходов конечно ($|\Omega| < \infty$) и все они равноправны (каждый выпадает с вероятностью $1 / |\Omega|$), задача сводится к чистой комбинаторике — доле «хороших» исходов среди всех:
\[
  P(A) = \frac{|A|}{|\Omega|}.
\]

\textbf{Условная вероятность и правило умножения.} Величина $P(A \mid B)$ означает вероятность наступления $A$, если мы \emph{уже знаем}, что событие $B$ точно произошло ($P(B) > 0$). Поскольку мир сузился до $B$, мы смотрим, какую долю внутри $B$ занимает пересечение $A \cap B$:
\[
  P(A \mid B) = \frac{P(A \cap B)}{P(B)} \quad \Longleftrightarrow \quad P(A \cap B) = P(B)\, P(A \mid B).
\]
Для цепочки шагов это даёт \textbf{теорему умножения} (при $P(A_1 \cap \dots \cap A_{k-1}) > 0$):
\[
  P(A_1 \cap A_2 \cap \dots \cap A_k) = P(A_1)\, P(A_2 \mid A_1)\, P(A_3 \mid A_1 \cap A_2) \dots P(A_k \mid A_1 \cap \dots \cap A_{k-1}).
\]

\textbf{Независимость событий.}
\begin{itemize}
  \item Два события $A$ и $B$ \textbf{независимы}, если информация о наступлении одного никак не меняет шансы другого ($P(A \mid B) = P(A)$), в универсальной записи без деления: $P(A \cap B) = P(A) P(B)$.
  \item Набор событий $A_1, \dots, A_n$ \textbf{независим в совокупности}, если для \emph{любого} непустого подмножества $I \subseteq [n]$ вероятность их одновременного наступления распадается в произведение: $P\bigl(\bigcap_{i \in I} A_i\bigr) = \prod_{i \in I} P(A_i)$. Важно: если каждые два события независимы попарно, весь набор целиком может не быть независимым в совокупности.
\end{itemize}

\textbf{Формула полной вероятности и формула Байеса.} Пусть всё пространство $\Omega$ разбито на непересекающиеся сценарии — \textbf{полную группу гипотез} $H_1, \dots, H_k$ ($\Omega = H_1 \sqcup \dots \sqcup H_k$, $P(H_i) > 0$), то есть в опыте гарантированно реализуется ровно одна из гипотез. Тогда:
\begin{itemize}
  \item \textbf{Формула полной вероятности} находит $P(A)$ разбором по тому, какой сценарий случился:
  \[
    P(A) = \sum_{i=1}^{k} P(H_i)\, P(A \mid H_i).
  \]
  \item \textbf{Формула Байеса} переворачивает условие — событие $A$ уже произошло ($P(A) > 0$), и мы переоцениваем вероятность того, что причиной была именно гипотеза $H_j$:
  \[
    P(H_j \mid A) = \frac{P(H_j)\, P(A \mid H_j)}{P(A)} = \frac{P(H_j)\, P(A \mid H_j)}{\sum_{i=1}^{k} P(H_i)\, P(A \mid H_i)}.
  \]
\end{itemize}

\textbf{Подсчёт вероятностей по простому модулю.} В олимпиадных задачах рациональный ответ — несократимую дробь $P / Q$ — просят вывести по простому модулю $p$ (обычно $10^9 + 7$ или $998244353$, при условии $p \nmid Q$) как число:
\[
  (P \bmod p) \cdot Q^{p-2} \bmod p.
\]
На практике \textbf{не нужно} хранить дроби в длинной арифметике до конца решения: складывать, вычитать, умножать и делить вероятности и матожидания можно сразу в остатках по модулю $p$ на каждом промежуточном шаге (в ДП, Гауссе, формулах), если ни на одном шаге не происходит деления на число, кратное $p$.

% ==============================================================================
\section{Уровень 2: Рабочая классика контестов}
% ==============================================================================

\subsection{Линейность математического ожидания}

\textbf{Случайная величина и её матожидание.}
\begin{itemize}
  \item \textbf{Дискретная случайная величина (СВ) $X$} — это просто числовая функция от исхода эксперимента: каждому исходу $\omega \in \Omega$ она ставит в соответствие число $X(\omega)$ (например, число очков, длина пути, количество инверсий). Её \emph{распределение} задаётся списком принимаемых значений $x$ и вероятностей $P(X = x)$ ($\sum_x P(X = x) = 1$).
  \item \textbf{Независимые СВ:} величины $X$ и $Y$ независимы, если при любых $x, y$ знание значения одной не влияет на шансы другой: $P(X = x \text{ и } Y = y) = P(X = x) P(Y = y)$.
  \item \textbf{Математическое ожидание ($\E[X]$):} среднее взвешенное значение величины $X$ (каждое возможное значение умножается на вероятность его получить; предполагается, что для бесконечных сумм ряд сходится абсолютно):
  \[
    \E[X] = \sum_{\omega \in \Omega} X(\omega) P(\{\omega\}) = \sum_{x} x \, P(X = x), \qquad
    \E[g(X)] = \sum_{x} g(x) \, P(X = x).
  \]
  \item \textbf{Дисперсия ($\Var(X)$):} числовая мера разброса СВ вокруг её среднего — средний квадрат отклонения, который на практике считают через $\E[X^2]$:
  \[
    \Var(X) = \E\bigl[(X - \E[X])^2\bigr] = \E[X^2] - (\E[X])^2.
  \]
\end{itemize}

\textbf{Линейность математического ожидания.} Среднее линейной комбинации равно такой же комбинации средних для \textbf{любых} случайных величин $X_1, \dots, X_n$ и констант $c_1, \dots, c_n \in \R$ — \textbf{даже если величины как угодно зависят друг от друга}:
\[
  \E\!\left[\sum_{i=1}^{n} c_i X_i\right] = \sum_{i=1}^{n} c_i \E[X_i].
\]
А вот мультипликативность $\E[XY] = \E[X]\E[Y]$ и сложение дисперсий $\Var(X + Y) = \Var(X) + \Var(Y)$ требуют независимости величин $X$ и $Y$.

\textbf{Метод индикаторов.} \textbf{Индикатор} события $A$ — простейшая случайная величина-флажок $\Ind_A$, равная $1$, если событие $A$ произошло, и $0$, если не произошло. У неё три рабочих свойства:
\[
  \E[\Ind_A] = 1 \cdot P(A) + 0 \cdot P(\bar{A}) = P(A), \qquad \Ind_A^2 = \Ind_A, \qquad \Ind_{A \cap B} = \Ind_A \cdot \Ind_B.
\]
Если нужно найти матожидание количества $X$ каких-то «хороших» объектов или сработавших событий из списка $A_1, \dots, A_m$ (число неподвижных точек, число компонент, число пар-инверсий), представляем $X = \sum_{i=1}^m \Ind_{A_i}$. По линейности зависимости между событиями игнорируются (а для $\E[X^2]$ достаточно попарных пересечений):
\[
  \E[X] = \sum_{i=1}^{m} P(A_i), \qquad
  \E[X^2] = \sum_{i=1}^{m} P(A_i) + 2 \sum_{1 \le i < j \le m} P(A_i \cap A_j).
\]

\textbf{Хвостовая формула (Tail Sum Formula).} Если случайная величина $X$ принимает только целые неотрицательные значения $\{0, 1, 2, \dots\}$, её матожидание можно считать через вероятности хвостов «$X$ хотя бы $k$» (часто $P(X \ge k)$ найти намного проще, чем $P(X = k)$, особенно когда $X$ — минимум или максимум):
\[
  \E[X] = \sum_{k=1}^{\infty} P(X \ge k) = \sum_{k=0}^{\infty} P(X > k), \qquad
  \E[X^2] = \sum_{k=1}^{\infty} (2k - 1) P(X \ge k).
\]

\textbf{Геометрическое распределение $T \sim \Geom(p)$.} Одиночный опыт с вероятностью успеха $p \in (0, 1]$ и неудачи $1 - p$ называется \emph{испытанием Бернулли}. Если повторять независимые испытания Бернулли подряд до первого успеха, то номер попытки $T \in \{1, 2, \dots\}$, на которой впервые случится успех, имеет геометрическое распределение:
\[
  P(T = k) = (1 - p)^{k-1} p \;\;(k \ge 1), \qquad P(T \ge k) = (1 - p)^{k-1}, \qquad
  \E[T] = \frac{1}{p}, \qquad \Var(T) = \frac{1 - p}{p^2}.
\]
\emph{Отсутствие памяти:} если первые $m$ попыток закончились неудачей, распределение оставшегося числа попыток точно такое же, как в самом начале: $P(T > m + k \mid T > m) = P(T > k)$.

\textbf{Задача о собирателе купонов (Coupon Collector's Problem).} Есть $n$ видов купонов, на каждом шаге равновероятно выпадает один случайный вид с шансом $1/n$. Сколько шагов $T$ в среднем нужно, чтобы собрать все $n$ видов? Разбиваем процесс на этапы: когда у нас уже есть $i$ различных видов, любой из оставшихся $n - i$ видов будет новым, то есть шанс поймать новый вид на одном шаге равен $p_i = \frac{n - i}{n}$, и мы ждём его $T_i \sim \Geom\bigl(\frac{n - i}{n}\bigr)$ шагов. Складывая ожидания этапов $T = T_0 + T_1 + \dots + T_{n-1}$:
\[
  \E[T] = \sum_{i=0}^{n-1} \frac{n}{n - i} = n H_n = n \sum_{k=1}^{n} \frac{1}{k} = n \ln n + \gamma n + O(1),
\]
где $H_n$ — $n$-е гармоническое число, $\gamma \approx 0{,}57721566$ — постоянная Эйлера—Маскерони. Чтобы собрать первые $m \le n$ различных видов, в среднем нужно $n(H_n - H_{n-m})$ шагов.

% ------------------------------------------------------------------------------
\subsection{Формула включений-исключений (ФВИ / PIE)}

\textbf{Суть формулы.} Пусть в множестве всех объектов $U$ выделены $k$ свойств (подмножеств) $A_1, \dots, A_k \subseteq U$. Чтобы посчитать объекты, обладающие \emph{хотя бы одним} свойством (объединение), или наоборот \emph{ни одним} из «запрещённых» свойств ($\bigcap_{i=1}^k \bar{A}_i$, где $\bar{A}_i = U \setminus A_i$), перебирают все подмножества условий $S \subseteq [k]$, для каждого считают число объектов, где выполнены сразу \textbf{все} условия из $S$, и складывают со знаком $(-1)^{|S|}$ (при $S = \emptyset$ пересечение равно всему $U$):
\[
  \left|\bigcup_{i=1}^{k} A_i\right| = \sum_{\emptyset \ne S \subseteq [k]} (-1)^{|S| - 1} \left|\bigcap_{i \in S} A_i\right|, \qquad
  \left|\bigcap_{i=1}^{k} \bar{A}_i\right| = \sum_{S \subseteq [k]} (-1)^{|S|} \left|\bigcap_{i \in S} A_i\right|.
\]
Точно так же для вероятностей: $P\bigl(\bigcup_{i=1}^k A_i\bigr) = \sum_{\emptyset \ne S \subseteq [k]} (-1)^{|S|-1} P\bigl(\bigcap_{i \in S} A_i\bigr)$. Если мощность пересечения $\bigl|\bigcap_{i \in S} A_i\bigr| = f(|S|)$ зависит только от числа выбранных условий $j = |S|$, то сумма из $2^k$ слагаемых сворачивается в $k + 1$ слагаемое: $\bigl|\bigcap_{i=1}^k \bar{A}_i\bigr| = \sum_{j=0}^k (-1)^j \binom{k}{j} f(j)$.

\textbf{Перебор масок за $O(2^k)$.} При небольшом числе условий ($k \le 20\text{--}25$) подмножество $S \subseteq [k]$ кодируется битовой маской $m \in [0, 2^k - 1]$, а знак слагаемого определяется чётностью числа единичных битов: $(-1)^{\popcnt(m)}$. Если пересчитывать пересечение рекурсивно или через отщепление младшего бита $b = m \mathbin{\&} (-m)$ за $O(1)$ от меньшей маски $m \oplus b$, весь перебор укладывается в чистые $O(2^k)$ вместо $O(k \cdot 2^k)$.

\textbf{Беспорядки (субфакториал $!n$).} \textbf{Беспорядок} (derangement) — перестановка $\pi \in S_n$ без неподвижных точек, в которой ни один элемент не остался на своём месте ($\pi_i \ne i$ для всех $i \in [n]$). По ФВИ (где $A_i$ — условие «$\pi_i = i$», а пересечение $j$ таких условий фиксирует $j$ позиций и даёт $(n - j)!$ вариантов на остальных) число беспорядков $!n$ равно:
\[
  !n = n! \sum_{j=0}^{n} \frac{(-1)^j}{j!}, \qquad !0 = 1, \quad !1 = 0, \quad !2 = 1.
\]
Удобные рекурренты для предподсчёта $!0, \dots, !N \bmod M$ за $O(N)$ по любому модулю $M$:
\[
  !n = (n - 1)\bigl(!(n - 1) + !(n - 2)\bigr) \;\; (n \ge 2), \qquad
  !n = n \cdot !(n - 1) + (-1)^n \;\; (n \ge 1).
\]
При $n \ge 1$: $!n = \bigl\lfloor \frac{n!}{e} + \frac{1}{2} \bigr\rfloor$, а доля беспорядков среди всех перестановок стремится к $\frac{!n}{n!} \to \frac{1}{e} \approx 0{,}3679$. Число перестановок в $S_n$ с ровно $m$ неподвижными точками равно $\binom{n}{m} \cdot !(n - m)$.

\textbf{ФВИ внутри ДП (чётность числа нарушенных условий).} Когда запрещённых условий много ($k \sim n$, например «между соседними элементами не должно быть запрещённой связи»), перебрать $2^k$ масок нельзя. Однако в формуле $\sum_S (-1)^{|S|} \left|\bigcap_{i \in S} A_i\right|$ нам не нужен сам список элементов $S$ — важен только знак $(-1)^{|S|}$ и то, как выбранные в $S$ нарушения ограничивают объект:
\begin{enumerate}
  \item \textbf{Чётность в состоянии ДП:} строим объект слева направо и храним в состоянии бит $p = |S| \bmod 2 \in \{0, 1\}$ — чётность числа нарушений, которые мы решили зафиксировать в $S$ (остальные переходы свободны — там условие может как нарушаться, так и нет). Когда на переходе мы берём очередное условие в $S$, меняем $p \leftarrow p \oplus 1$. Ответ в конце: $\mathrm{dp}_{\text{end}}[0] - \mathrm{dp}_{\text{end}}[1]$.
  \item \textbf{Вшивание знака $(-1)$ сразу в переход:} два состояния $p = 0$ и $p = 1$ можно сразу объединить, храня в ДП сразу разность «чётные минус нечётные»: тогда выбор очередного нарушенного условия просто умножает коэффициент перехода на $(-1)$ (а склейка блока длины $L$ из $L - 1$ нарушений подряд даёт множитель $(-1)^{L - 1}$).
\end{enumerate}

% ------------------------------------------------------------------------------
\subsection{Числа Каталана и принцип отражений (Андре)}

\textbf{Числа Каталана ($C_n$, $n \ge 0$).} Последовательность $(C_n)_{n \ge 0} = (1, 1, 2, 5, 14, 42, 132, \ldots)$:
\[
  C_n = \frac{1}{n + 1}\binom{2n}{n} = \binom{2n}{n} - \binom{2n}{n + 1} \quad (n \ge 0).
\]
Рекурренты (квадратичная за $O(n^2)$ работает без деления по любому модулю $M$ и отвечает разбиению объекта на «первый неделимый кусок» размера $i$ и «остаток» размера $n - 1 - i$):
\[
  C_0 = 1, \qquad C_n = \sum_{i=0}^{n-1} C_i C_{n - 1 - i} \;\;(n \ge 1), \qquad
  C_n = \frac{2(2n - 1)}{n + 1}\, C_{n - 1}.
\]

\textbf{Главные комбинаторные объекты, количество которых равно $C_n$:}
\begin{enumerate}
  \item \textbf{Правильные скобочные последовательности (ПСП)} из $n$ открывающих и $n$ закрывающих скобок (длина $2n$): если считать \texttt{(} за $+1$, а \texttt{)} за $-1$, то сумма на любом префиксе $\ge 0$, а общая сумма равна $0$.
  \item \textbf{Пути на сетке не выше диагонали (пути Дика):} пути из $(0, 0)$ в $(n, n)$ шагами вправо $(1, 0)$ и вверх $(0, 1)$, ни в одной точке не заходящие выше прямой $y = x$ ($y \le x$ везде). Если повернуть картинку на $45^\circ$, это пути из $(0, 0)$ в $(2n, 0)$ шагами $(1, +1)$ и $(1, -1)$, не опускающиеся ниже оси $y = 0$.
  \item \textbf{Бинарные и плоские корневые деревья:} корневые бинарные деревья (где различаются левый и правый сын) с $n$ вершинами; полные бинарные деревья с $n + 1$ листьями; произвольные плоские корневые деревья (где важен порядок сыновей) на $n + 1$ вершинах.
  \item \textbf{Триангуляции выпуклого $(n + 2)$-угольника:} способы разрезать выпуклый $(n + 2)$-угольник с фиксированными вершинами на $n$ треугольников с помощью $n - 1$ непересекающихся диагоналей.
  \item \textbf{Хорды и стековые перестановки:} способы соединить $2n$ точек на окружности $n$ непересекающимися хордами; перестановки длины $n$, которые можно отсортировать одним стеком (не содержат подпоследовательности вида $2\text{--}3\text{--}1$).
\end{enumerate}

\textbf{Принцип отражений Андре.} Без ограничений число путей по сетке из $(x_0, y_0)$ в $(x_1, y_1)$ шагами $(1, 0)$ и $(0, 1)$ равно $\binom{\Delta x + \Delta y}{\Delta x}$, где $\Delta x = x_1 - x_0 \ge 0$, $\Delta y = y_1 - y_0 \ge 0$ (выбираем, какие $\Delta x$ шагов из $\Delta x + \Delta y$ будут вправо). Пусть теперь проведена запретная прямая $L : y = x + c$ ($c \in \Z$), которую пути \textbf{нельзя даже касаться}, и обе точки лежат строго ниже неё ($y_0 < x_0 + c$, $y_1 < x_1 + c$; если в задаче разрешено касаться прямой $y = x + d$, но нельзя подниматься выше, ставим запретную прямую на единицу выше: $c = d + 1$).

Любой «плохой» путь, коснувшийся $L$, отразим относительно $L$ на участке от старта до \emph{первой} точки касания — старт $(x_0, y_0)$ перейдёт в зеркальную точку $(y_0 - c, x_0 + c)$ по другую сторону $L$, и получится взаимно однозначное соответствие между всеми плохими путями и вообще всеми путями из $(y_0 - c, x_0 + c)$ в $(x_1, y_1)$. Вычитая плохие пути из всех, получаем число допустимых путей строго ниже $L$:
\[
  N_{< L} = \binom{\Delta x + \Delta y}{\Delta x} - \binom{\Delta x + \Delta y}{\Delta x + c - (y_0 - x_0)}.
\]
\emph{Частный случай:} число путей из $(0, 0)$ в $(n, m)$ ($n \ge m \ge 0$), не поднимающихся выше $y = x$ ($c = 1$), равно $\binom{n + m}{n} - \binom{n + m}{n + 1} = \frac{n - m + 1}{n + 1}\binom{n + m}{m}$ (при $m = n$ даёт $C_n$).

\textbf{Обобщение: пути строго между двумя параллельными прямыми.} Пусть пути из $(x_0, y_0)$ в $(x_1, y_1)$ запрещено касаться сразу двух прямых — верхней $L_1 : y = x + a$ и нижней $L_2 : y = x - b$ ($a, b \in \Z$, обе точки лежат строго в коридоре: $x_i - b < y_i < x_i + a$; для нестрогого коридора $x - b' \le y \le x + a'$ берём $a = a' + 1$, $b = b' + 1$). Пусть $D = a + b \ge 2$ — ширина коридора. Так как путь может по очереди биться то о верхнюю, то о нижнюю границу, одно отражение не спасает — делаем цепочку последовательных отражений старта от $L_1$ и $L_2$ с чередованием знаков по ФВИ:
\[
  N_{L_2 < \text{путь} < L_1} = \sum_{k \in \Z} \binom{\Delta x + \Delta y}{\Delta x - kD} - \sum_{k \in \Z} \binom{\Delta x + \Delta y}{\Delta x + a - (y_0 - x_0) + kD}, \qquad D = a + b.
\]
Ненулевыми здесь будут только $O\bigl(\frac{\Delta x + \Delta y}{D}\bigr)$ слагаемых, где нижний индекс попадает в $[0, \Delta x + \Delta y]$.

% ------------------------------------------------------------------------------
\subsection{ДП по вероятностям и матожиданиям}

\textbf{Модель состояний и переходов.} Процесс стартует в состоянии $s \in V$ и заканчивается, как только попадает в одно из терминальных (конечных) состояний $T \subseteq V$. Из каждого нетерминального состояния $u \notin T$ за один шаг происходит случайный переход в $v$ с вероятностью $p_{u,v} \ge 0$ ($\sum_{v \in V} p_{u,v} = 1$) и стоимостью $c_{u,v}$ (если считаем число шагов, то $c_{u,v} = 1$). Обозначим среднюю цену одного шага из $u$ через $\bar{c}_u = \sum_{v \in V} p_{u,v} c_{u,v}$ (для числа шагов $\bar{c}_u = 1$). Предполагаем пока, что граф переходов не содержит циклов длины $\ge 2$ (DAG).

\textbf{Прямое ДП (от старта к финишу) — естественно для вероятностей.} Обозначим через $P[u]$ вероятность того, что процесс вообще когда-нибудь побывает в состоянии $u$. Идём в топологическом порядке от $s$:
\[
  P[s] = 1, \qquad P[v] = \sum_{u : u \to v} P[u] \cdot p_{u,v}.
\]
\emph{Важная ловушка для матожиданий:} писать прямое ДП вида $E[v] = \sum_{u \to v} p_{u,v}(E[u] + c_{u,v})$ \textbf{нельзя}, потому что $p_{u,v}$ — это шанс пойти в $v$, уже стоя в $u$, а в сумме они даже не дают $1$ по входящим рёбрам! Если по какой-то причине нужен именно прямой проход для среднего пути до $v$, хранят пару $(P[u], S[u])$, где $S[u]$ — суммарная стоимость путей до $u$, взвешенная их вероятностями:
\[
  S[s] = 0, \qquad S[v] = \sum_{u : u \to v} \bigl(S[u] + P[u] \cdot c_{u,v}\bigr) p_{u,v}, \qquad
  \E[\text{стоимость до } v \mid \text{мы попали в } v] = \frac{S[v]}{P[v]}.
\]

\textbf{Обратное ДП (от финиша к старту) — золотой стандарт для матожиданий.} Заведём $E[u]$ — сколько \emph{ещё в среднем осталось} заплатить до конца игры, если сейчас мы стоим в состоянии $u$ (и аналогично $f[u]$ — вероятность выиграть, попав в $T_{\text{win}} \subseteq T$, стартуя из $u$). В этой постановке сумма $p_{u,v}$ по исходящим переходам равна $1$, никаких условных делений нет, пересчёт идёт в обратном топологическом порядке, а ответ задачи — просто $E[s]$ (или $f[s]$):
\[
  E[t] = 0 \;\; (t \in T), \qquad E[u] = \bar{c}_u + \sum_{v \in V} p_{u,v} E[v] \;\; (u \notin T);
\]
\[
  f[t] = \iverson{t \in T_{\text{win}}} \;\; (t \in T), \qquad f[u] = \sum_{v \in V} p_{u,v} f[v] \;\; (u \notin T).
\]

\textbf{Избавление от петель (переход из состояния в себя).} Если из состояния $u$ с вероятностью $p_{u,u} \in (0, 1)$ ход оказывается «холостым» и оставляет нас в том же $u$, в правой части уравнения появляется $p_{u,u} E[u]$. Переносим его влево, приводим подобные и делим на $(1 - p_{u,u})$ — вероятность с очередной попытки покинуть $u$:
\[
  E[u] = \frac{\bar{c}_u + \sum_{v \ne u} p_{u,v} E[v]}{1 - p_{u,u}}, \qquad
  f[u] = \frac{\sum_{v \ne u} p_{u,v} f[v]}{1 - p_{u,u}}.
\]
Для числа шагов ($\bar{c}_u = 1$) числитель $E[u]$ равен $1 + \sum_{v \ne u} p_{u,v} E[v]$ (не забываем, что холостые шаги в петлю тоже тратят $1$ ход и уже учтены в $\bar{c}_u = 1$!).

% ==============================================================================
\section{Уровень 3: Комбинаторные числа, графы и СЛАУ}
% ==============================================================================

\subsection{Числа Стирлинга и числа Белла}

\textbf{Числа Стирлинга 1-го рода (перестановки с $k$ циклами).} \textbf{Беззнаковое число Стирлинга 1-го рода} $\stirlingfirst{n}{k}$ (или $c(n, k)$, $0 \le k \le n$) — это количество перестановок длины $n$, которые распадаются ровно на $k$ независимых циклов (наглядно: сколькими способами рассадить $n$ человек за $k$ неразличимых круглых столов без пустых столов). Рекуррента за $O(n^2)$ (новый $n$-й человек либо садится за отдельный стол, либо встаёт слева от любого из $n - 1$ уже сидящих):
\[
  \stirlingfirst{0}{0} = 1, \qquad \stirlingfirst{n}{0} = \stirlingfirst{0}{k} = 0 \;\;(n, k \ge 1), \qquad
  \stirlingfirst{n}{k} = (n - 1)\stirlingfirst{n - 1}{k} + \stirlingfirst{n - 1}{k - 1}.
\]
\textbf{Знаковые числа Стирлинга 1-го рода} $s(n, k) = (-1)^{n - k}\stirlingfirst{n}{k}$ появляются при раскрытии скобок в убывающем факториале (а беззнаковые — в возрастающем):
\[
  \rise{x}{n} = \sum_{k=0}^{n} \stirlingfirst{n}{k} x^k, \qquad
  \fall{x}{n} = \sum_{k=0}^{n} (-1)^{n - k} \stirlingfirst{n}{k} x^k, \qquad
  \sum_{k=0}^{n} \stirlingfirst{n}{k} = n!.
\]

\textbf{Числа Стирлинга 2-го рода (разбиение множества на $k$ кучек).} Число $\stirlingsecond{n}{k}$ (или $S(n, k)$, $0 \le k \le n$) — количество способов разбить $n$ различных элементов на $k$ непустых \textbf{ненумерованных} кучек (блоков). Если кучки пронумеровать, получится $k! \stirlingsecond{n}{k}$ — число сюръекций (отображений из $[n]$ в $[k]$ без пустых ящиков). Рекуррента за $O(n^2)$ ($n$-й элемент либо образует новую кучку, либо идёт в любую из $k$ существующих):
\[
  \stirlingsecond{0}{0} = 1, \qquad \stirlingsecond{n}{0} = \stirlingsecond{0}{k} = 0 \;\;(n, k \ge 1), \qquad
  \stirlingsecond{n}{k} = k\stirlingsecond{n - 1}{k} + \stirlingsecond{n - 1}{k - 1}.
\]
Явная формула через ФВИ по пустым ящикам (считает одно число за $O(k \log n)$ или всю строку $\stirlingsecond{n}{0 \dots n}$ одной свёрткой NTT за $O(n \log n)$):
\[
  \stirlingsecond{n}{k} = \frac{1}{k!} \sum_{j=0}^{k} (-1)^{k - j} \binom{k}{j} j^n = \sum_{j=0}^{k} \frac{(-1)^{k - j}}{(k - j)!} \cdot \frac{j^n}{j!}.
\]

\textbf{Числа Белла ($B_n$).} Общее количество способов разбить $n$ различных элементов на \emph{любое} число непустых ненумерованных кучек ($B_0 = 1, B_1 = 1, B_2 = 2, B_3 = 5, B_4 = 15$; полезно для оценки полного перебора разбиений: $B_{10} = 115\,975$, $B_{12} = 4\,213\,597$):
\[
  B_n = \sum_{k=0}^{n} \stirlingsecond{n}{k}, \qquad
  B_{n + 1} = \sum_{k=0}^{n} \binom{n}{k} B_k, \qquad
  \sum_{n=0}^{\infty} B_n \frac{x^n}{n!} = \exp(e^x - 1).
\]
По простому модулю $p$ выполняется сравнение Ташарда: $B_{n + p} \equiv B_n + B_{n + 1} \pmod p$.

\textbf{Переход от $x^k$ к убывающим факториалам и подсчёт сумм $\sum_i f(i)^k$.} Обычная степень $x^k$ раскладывается по убывающим факториалам или биномиальным коэффициентам с коэффициентами Стирлинга 2-го рода (верно при любых $x \in \R$ и целых $k \ge 0$):
\[
  x^k = \sum_{j=0}^{k} \stirlingsecond{k}{j} \fall{x}{j} = \sum_{j=0}^{k} \stirlingsecond{k}{j} j! \binom{x}{j}.
\]
Зачем это нужно на контестах: когда $k$ небольшое ($k \le 500$), а слагаемых очень много, со степенью $f(i)^k$ неудобно работать напрямую, зато с биномами $\binom{f(i)}{j}$ работают все классические тождества:
\begin{itemize}
  \item \textbf{Сумма $k$-х степеней $1^k + 2^k + \dots + n^k$} за $O(k^2)$ (подставляем разложение и сворачиваем внутреннюю сумму по $i$ хоккейной клюшкой):
  \[
    \sum_{i=0}^{n} i^k = \sum_{j=0}^{k} \stirlingsecond{k}{j} j! \binom{n + 1}{j + 1} = \sum_{j=0}^{k} \stirlingsecond{k}{j} \frac{\fall{(n + 1)}{j + 1}}{j + 1}.
  \]
  \item \textbf{Сумма вида $\sum_{i=0}^n \binom{n}{i} i^k$} за $O(k)$:
  \[
    \sum_{i=0}^{n} \binom{n}{i} i^k = \sum_{j=0}^{k} \stirlingsecond{k}{j} \fall{n}{j} 2^{n - j}.
  \]
  \item \textbf{ДП для сумм $\sum_{\omega} f(\omega)^k$ (например, сумма $k$-х степеней длин всех путей в дереве):} вместо степеней храним в вершине массив $D_j = \sum_{\omega} \binom{f(\omega)}{j}$ для $0 \le j \le k$. При удлинении всех путей на $1$ ($f \to f + 1$) массив пересчитывается по Паскалю $D'_j = D_j + D_{j - 1}$ за $O(k)$, при склейке двух поддеревьев ($f_1 + f_2$) — свёрткой Вандермонда $D_j = \sum_{t=0}^j D^{(1)}_t D^{(2)}_{j - t}$ за $O(k^2)$, а в конце ответ собирается как $\sum_{j=0}^k \stirlingsecond{k}{j} j! \, D_j$.
\end{itemize}

% ------------------------------------------------------------------------------
\subsection{Теорема Люка и китайская теорема об остатках (КТО / CRT)}

\textbf{Теорема Люка (маленький простой модуль $p$, огромное $n \le 10^{18}$).} Когда $n \ge p$, в $n!$ входит множитель $p$, поэтому формула через \texttt{inv\_fact[n]} перестаёт работать. Теорема Люка говорит: запишем $n$ и $k$ в системе счисления с основанием $p$: $n = \sum_{i=0}^s n_i p^i$, $k = \sum_{i=0}^s k_i p^i$ ($0 \le n_i, k_i < p$). Тогда биномиальный коэффициент по модулю $p$ просто распадается в произведение по разрядам:
\[
  \binom{n}{k} \equiv \prod_{i=0}^{s} \binom{n_i}{k_i} \pmod p.
\]
Предподсчитав факториалы до $p - 1$ за $O(p)$, отвечаем на запрос за $O(\log_p n)$. Полезные следствия: $\binom{n}{k}$ кратно $p$ тогда и только тогда, когда хотя бы в одном разряде $k_i > n_i$; при $p = 2$: $\binom{n}{k}$ нечётно $\iff (k \mathbin{\&} n) = k$ (единичные биты $k$ — подмаска битов $n$).

\textbf{Китайская теорема об остатках (КТО / CRT).} Если известны остатки неизвестного числа $x$ по попарно взаимно простым модулям $m_1, \dots, m_r$ ($\gcd(m_i, m_j) = 1$ при $i \ne j$), то остаток $x$ по их произведению $M = m_1 \dots m_r$ восстанавливается однозначно:
\[
  x \equiv \sum_{i=1}^{r} a_i M_i y_i \pmod M, \qquad M_i = \frac{M}{m_i}, \qquad y_i = M_i^{-1} \bmod m_i.
\]
\emph{Если модули не взаимно просты:} два сравнения $x \equiv a_1 \pmod{m_1}$ и $x \equiv a_2 \pmod{m_2}$ не противоречат друг другу тогда и только тогда, когда $a_1 \equiv a_2 \pmod g$, где $g = \gcd(m_1, m_2)$. В этом случае они склеиваются в одно сравнение по модулю $\lcm(m_1, m_2)$: находим $\frac{m_1}{g} u + \frac{m_2}{g} v = 1$ расширенным Евклидом и берём $x = a_1 + m_1 \cdot \bigl(\frac{a_2 - a_1}{g} u \bmod \frac{m_2}{g}\bigr)$.

\textbf{Как считать $\binom{n}{k} \bmod M$ при $n \le 10^{18}$ и составном модуле $M = \prod_{i=1}^r p_i^{\alpha_i}$:}
\begin{enumerate}
  \item \textbf{$M$ свободно от квадратов ($M = p_1 \dots p_r$, все $\alpha_i = 1$):} находим остаток $\binom{n}{k} \bmod p_i$ по теореме Люка для каждого простого $p_i$ и собираем ответ через КТО.
  \item \textbf{В $M$ есть степени простых $p^\alpha$ ($\alpha \ge 2$, метод Грэнвилла):} для каждой степени $p^\alpha$ считаем $\binom{n}{k} \bmod p^\alpha$ отдельно, отделяя множители $p$ от взаимно простой с $p$ части факториала:
  \begin{itemize}
    \item \emph{Степень вхождения $p$:} по \textbf{формуле Лежандра} простое $p$ входит в $m!$ в степени $v_p(m!) = \sum_{j \ge 1} \lfloor m / p^j \rfloor$. Значит, в $\binom{n}{k}$ оно входит в степени $e = v_p(n!) - v_p(k!) - v_p((n - k)!)$ (по \textbf{теореме Куммера} $e$ также равно числу переносов через разряд при сложении $k + (n - k)$ столбиком в базе $p$). Если $e \ge \alpha$, остаток сразу равен $0$.
    \item \emph{Оставшаяся часть без множителей $p$:} представим $m! = p^{v_p(m!)} f_p(m)$, где $\gcd(f_p(m), p) = 1$. Если в произведении $1 \cdot 2 \dots m$ вычеркнуть числа, кратные $p$, оставшееся произведение $(m!)_p = \prod_{1 \le j \le m,\, p \nmid j} j \bmod p^\alpha$ периодично с длиной периода $p^\alpha$:
    \[
      (m!)_p \equiv U^{\lfloor m / p^\alpha \rfloor} \cdot \bigl((m \bmod p^\alpha)!\bigr)_p \pmod{p^\alpha}, \qquad
      U = (p^\alpha!)_p \equiv \begin{cases} 1 \pmod{p^\alpha}, & p = 2,\; \alpha \ge 3, \\ -1 \pmod{p^\alpha}, & \text{иначе.} \end{cases}
    \]
    Вынесенные же кратные $p$ после сокращения на $p$ превращаются в $1 \cdot 2 \dots \lfloor m / p \rfloor$, поэтому вся взаимно простая часть рекурсивно собирается за $O(\log_p m)$ (после предподсчёта префиксных произведений длины $p^\alpha$ за $O(p^\alpha)$): $f_p(m) \equiv \prod_{j \ge 0} \bigl(\lfloor m / p^j \rfloor !\bigr)_p \pmod{p^\alpha}$.
    \item Так как $f_p(k)$ и $f_p(n - k)$ взаимно просты с $p^\alpha$, их можно обратить расширенным Евклидом:
    \[
      \binom{n}{k} \equiv p^e \cdot f_p(n) \cdot \bigl(f_p(k)\bigr)^{-1} \cdot \bigl(f_p(n - k)\bigr)^{-1} \pmod{p^\alpha},
    \]
    а затем склеить ответы по всем $p_i^{\alpha_i}$ через КТО.
  \end{itemize}
\end{enumerate}

% ------------------------------------------------------------------------------
\subsection{Цепи Маркова и метод Гаусса}

\textbf{Цепь Маркова — случайное блуждание без памяти о прошлом.} Пусть система шагами переходит между конечным числом состояний $V = \{1, \dots, N\}$ ($X_t$ — состояние на шаге $t$). Процесс называется (однородной) \textbf{цепью Маркова}, если вероятности следующего шага зависят \emph{только от текущего состояния}, а не от того, каким маршрутом мы в него попали (\emph{марковское свойство}):
\[
  P(X_{t + 1} = j \mid X_t = i, X_{t - 1} = i_{t - 1}, \dots, X_0 = i_0) = P(X_{t + 1} = j \mid X_t = i) = p_{i,j}.
\]
\begin{itemize}
  \item \textbf{Матрица переходов $P = (p_{i,j})_{N \times N}$:} сумма в каждой строке равна $1$ ($\sum_{j=1}^N p_{i,j} = 1$). Если вектор-строка $\pi^{(0)}$ хранит вероятности находиться в состояниях в начале, то за один шаг распределение умножается на $P$, а через $t$ шагов равно $\pi^{(t)} = \pi^{(0)} P^t$ (при $t \le 10^{18}$ и малом $N$ считается бинарным возведением матрицы в степень за $O(N^3 \log t)$).
  \item \textbf{Поглощающие и проходные состояния:} состояние $t$ \emph{поглощающее}, если из него нет выхода ($p_{t,t} = 1$ — финиш игры); состояние \emph{невозвратное (проходное)}, если процесс с ненулевой вероятностью покидает его навсегда.
  \item \textbf{Стационарное распределение:} вектор вероятностей $\pi = (\pi_1, \dots, \pi_N)$ ($\sum \pi_i = 1, \pi_i \ge 0$), который не меняется после шага цепи: $\pi P = \pi$ (находится Гауссом). Если из любого состояния можно добраться в любое другое и цепь не зацикливается строго по чётности шагов, распределение $\pi^{(t)}$ при $t \to \infty$ сходится к единственному $\pi$, а $\pi_i$ показывает долю времени, проводимую в состоянии $i$ (обратную к среднему числу шагов до первого возвращения в $i$).
\end{itemize}

\textbf{Циклическое ДП по матожиданиям через метод Гаусса за $O(N^3)$.} Если в графе переходов между нетерминальными состояниями $Q = V \setminus T$ ($|Q| = m$) есть циклы длины $\ge 2$, посчитать $E[u]$ в топологическом порядке нельзя. Зато уравнения обратного ДП ($E[t] = 0$ для $t \in T$) для всех $u \in Q$ линейны:
\[
  E[u] - \sum_{v \in Q} p_{u,v} E[v] = \bar{c}_u \quad (u \in Q) \quad \Longleftrightarrow \quad (I - P_{QQ}) \mathbf{E} = \bar{\mathbf{c}}.
\]
Если из каждого состояния $u \in Q$ есть ненулевой шанс рано или поздно дойти до терминального множества $T$, эта СЛАУ гарантированно имеет единственное решение и решается методом Гаусса за $O(m^3)$ (в \texttt{double} с выбором максимального по модулю опорного элемента в столбце, либо по модулю $p$ через обратные элементы). При этом элемент обратной матрицы $F = (I - P_{QQ})^{-1} = \sum_{k=0}^\infty P_{QQ}^k$ в клетке $(u, v)$ равен среднему числу раз, которое мы побываем в $v$ до конца игры при старте из $u$.

\textbf{Ускорение Гаусса на специальных графах:}
\begin{itemize}
  \item \textbf{Ленточные матрицы за $O(N W^2)$:} если состояния можно пронумеровать так, что переходы идут только в номера на расстоянии $\le W$ (то есть $A_{i,j} = 0$ при $|i - j| > W$; классический пример — таблица $H \times W$ при $W \ll H$, пронумерованная построчно, где сосед снизу отстоит на $+W$), то при занулении $i$-го столбца методом Гаусса нужно обработать всего $W$ строк ниже, меняя в каждой не более $2W$ ненулевых элементов. Асимптотика падает с $O(N^3)$ до $O(N W^2)$. При $W = 1$ (блуждание по полоске $1 \times N$, трёхдиагональная матрица) это даёт \textbf{метод прогонки} за $O(N)$.
  \item \textbf{Гаусс на дереве за $O(N)$ (прогонка по поддеревьям):} пусть переходы идут только по рёбрам дерева (из вершины $u$ — в родителя $p_u$ и в сыновей $v \in \operatorname{ch}(u)$). Подвесим дерево за корень $r$. Идея: для любого поддерева внешний мир виден только через родителя $p_u$, поэтому при обходе снизу вверх (DFS post-order) мы можем линейно выразить $E[u]$ через одно неизвестное $E[p_u]$:
  \[
    E[u] = k_u \cdot E[p_u] + b_u.
  \]
  В терминальных вершинах $u \in T$ сразу $k_u = 0, b_u = 0$. В нетерминальной вершине $u \notin T$ у всех её сыновей $v \in \operatorname{ch}(u)$ коэффициенты $(k_v, b_v)$ уже посчитаны; подставляем $E[v] = k_v E[u] + b_v$ в уравнение вершины $E[u] = \bar{c}_u + p_{u, p_u} E[p_u] + \sum_v p_{u,v} E[v]$, собираем $E[u]$ слева и выражаем:
  \[
    \Delta_u = 1 - \sum_{v \in \operatorname{ch}(u)} p_{u,v} k_v, \qquad
    k_u = \frac{p_{u, p_u}}{\Delta_u}, \qquad
    b_u = \frac{\bar{c}_u + \sum_{v \in \operatorname{ch}(u)} p_{u,v} b_v}{\Delta_u}.
  \]
  У корня $r$ родителя нет ($p_{r, p_r} = 0$), поэтому на самом верху сразу получаем число $E[r] = b_r$, а при необходимости одним проходом сверху вниз находим все остальные $E[u] = k_u E[p_u] + b_u$ за $O(N)$.
\end{itemize}

% ------------------------------------------------------------------------------
\subsection{Лемма Бёрнсайда и теорема Пойа}

\textbf{Симметрии, орбиты и неподвижные точки.} В задачах вида «посчитать раскраски с точностью до поворотов и отражений» у нас есть множество всех обычных фиксированных объектов $X$ и набор допустимых преобразований симметрии $G$ (включая холостое преобразование «ничего не поворачивать»; композиция двух симметрий и обратная симметрия тоже лежат в $G$ — такое семейство образует \emph{группу симметрий}):
\begin{itemize}
  \item \textbf{Орбита} объекта $x \in X$ ($G \cdot x = \{g \cdot x : g \in G\}$) — множество всех его повёрнутых/отражённых копий. Две раскраски считаются одинаковыми тогда и только тогда, когда лежат в одной орбите, поэтому \textbf{искомое число различных объектов — это в точности число орбит} $|X / G|$.
  \item \textbf{Стабилизатор} $\Stab_G(x) = \{g \in G : g \cdot x = x\}$ — все симметрии, при которых сам объект $x$ переходит строго в себя (всегда $|G \cdot x| \cdot |\Stab_G(x)| = |G|$: чем симметричнее объект, тем меньше у него различных копий в орбите).
  \item \textbf{Неподвижные точки симметрии $g \in G$:} множество $\Fix(g) = \{x \in X : g \cdot x = x\}$ всех объектов, которые вообще не меняются, если к ним применить конкретное преобразование $g$.
\end{itemize}

\textbf{Лемма Бёрнсайда.} Чтобы не отсеивать дубликаты внутри каждой орбиты, задача переворачивается: перебираем все симметрии $g \in G$, для каждой считаем, сколько объектов не меняются при применении $g$, и берём среднее арифметическое:
\[
  |X / G| = \frac{1}{|G|} \sum_{g \in G} |\Fix(g)|.
\]

\textbf{Теорема Пойа (для раскрасок).} Пусть мы красим $n$ позиций (вершин, клеток, граней) в $m$ цветов, а каждая симметрия $g \in G$ просто переставляет эти $n$ позиций местами. Как любая перестановка, $g$ распадается на $c(g)$ независимых циклов (обозначим через $c_j(g)$ количество циклов длины $j$).
\begin{itemize}
  \item \textbf{Без ограничений на число цветов:} чтобы раскраска не изменилась при сдвиге по циклам $g$, все клетки внутри каждого цикла обязаны быть покрашены в один цвет. Выбирая цвет независимо для каждого из $c(g)$ циклов, получаем $|\Fix(g)| = m^{c(g)}$ и готовую формулу:
  \[
    \text{Число различных раскрасок} = \frac{1}{|G|} \sum_{g \in G} m^{c(g)}.
  \]
  \item \textbf{Весовая форма (фиксировано, сколько клеток какого цвета нужно взять):} выписываем \emph{цикловой индекс} $Z_G(t_1, \dots, t_n) = \frac{1}{|G|} \sum_{g \in G} \prod_{j=1}^n t_j^{c_j(g)}$ и вместо каждой переменной $t_j$ (отвечающей циклу длины $j$, который целиком забирает $j$ клеток одного цвета) подставляем сумму $w_1^j + w_2^j + \dots + w_m^j$:
  \[
    F(w_1, \dots, w_m) = Z_G\!\left(\sum_{r=1}^{m} w_r,\; \sum_{r=1}^{m} w_r^2,\; \dots,\; \sum_{r=1}^{m} w_r^n\right).
  \]
  Коэффициент при одночлене $w_1^{k_1} \dots w_m^{k_m}$ равен числу неэквивалентных раскрасок, где $1$-й цвет использован $k_1$ раз, $\dots$, $m$-й цвет — $k_m$ раз.
\end{itemize}

\textbf{Готовые формулы для круга из $n$ бусин в $m$ цветов:}
\begin{itemize}
  \item \textbf{Ожерелья (только $n$ поворотов, группа $C_n$):} поворот на $k \in [0, n - 1]$ позиций распадается на $\gcd(n, k)$ циклов одинаковой длины $n / \gcd(n, k)$. Группируя одинаковые значения $d = \gcd(n, k) \mid n$ (таких $k$ ровно $\varphi(n / d)$ штук), получаем расчёт за $O(\sqrt{n})$ по делителям $n$:
  \[
    N_{\mathrm{rot}}(n, m) = \frac{1}{n} \sum_{k=0}^{n-1} m^{\gcd(n, k)} = \frac{1}{n} \sum_{d \mid n} \varphi\!\left(\frac{n}{d}\right) m^d.
  \]
  \item \textbf{Браслеты ($n$ поворотов + $n$ переворачиваний/отражений, группа $D_n$, $|D_n| = 2n$):} при нечётном $n$ ось каждого отражения проходит через $1$ бусину и середину противоположного ребра ($\frac{n + 1}{2}$ циклов); при чётном $n$ половина осей проходит через $2$ противоположные бусины ($\frac{n}{2} + 1$ циклов), а половина — между бусинами ($\frac{n}{2}$ циклов):
  \[
    N_{\mathrm{dih}}(n, m) =
    \begin{cases}
      \dfrac{1}{2n}\!\left(\displaystyle\sum_{d \mid n} \varphi\!\left(\frac{n}{d}\right) m^d + n \cdot m^{(n + 1)/2}\right), & n \text{ нечётно}, \\[10pt]
      \dfrac{1}{2n}\!\left(\displaystyle\sum_{d \mid n} \varphi\!\left(\frac{n}{d}\right) m^d + \frac{n}{2} m^{n/2 + 1} + \frac{n}{2} m^{n/2}\right), & n \text{ чётно.}
    \end{cases}
  \]
\end{itemize}

% ------------------------------------------------------------------------------
\subsection{Комбинаторика на графах}

\textbf{Код Прюфера.} Взаимно однозначный способ закодировать любое дерево на пронумерованных вершинах $\{1, \dots, n\}$ ($n \ge 2$) массивом из $n - 2$ чисел $(p_1, \dots, p_{n - 2})$, где каждое $p_i \in \{1, \dots, n\}$:
\begin{itemize}
  \item \textbf{Дерево $\to$ код:} $n - 2$ раза находим лист с \emph{минимальным} номером, выписываем в код его единственного соседа и удаляем этот лист из дерева (в конце останется одно ребро).
  \item \textbf{Код $\to$ дерево и главное свойство степеней:} каждая вершина $v$ встречается в коде Прюфера ровно $\deg(v) - 1$ раз (в частности, листья — те и только те вершины, которых нет в коде). Поэтому изначально зная все степени $\deg(v) = 1 + \mathrm{cnt}_p(v)$, мы $n - 2$ раза берём минимальную вершину со степенью $1$, соединяем её ребром с очередным $p_i$ и уменьшаем степени обеих вершин на $1$, а в финале соединяем две последние вершины со степенью $1$. С поддержкой указателя на минимальный лист оба перевода работают линейно за $O(n)$.
\end{itemize}

\textbf{Формула Кэли и подсчёт деревьев по коду Прюфера:}
\begin{itemize}
  \item \textbf{Формула Кэли:} так как массив длины $n - 2$ из чисел $1 \dots n$ может быть любым, число всех деревьев на $n$ занумерованных вершинах равно $n^{n - 2}$ (при $n = 1$ равно $1$).
  \item \textbf{Фиксированные степени вершин:} число деревьев на $n \ge 2$ вершинах с заданными степенями $d_1, \dots, d_n \ge 1$ ($d_1 + \dots + d_n = 2n - 2$) равно числу перестановок кода Прюфера с кратностями $d_i - 1$: $\frac{(n - 2)!}{(d_1 - 1)!\,(d_2 - 1)!\dots(d_n - 1)!}$.
  \item \textbf{Сшивание компонент связности:} если граф на $n$ вершинах уже разбит на $k \ge 1$ связных компонент размеров $s_1, \dots, s_k$, то число способов добавить $k - 1$ ребро, чтобы сделать граф связным, равно $n^{k - 2} \prod_{i=1}^k s_i$ (при $k = 1$ равно $1$). Частный случай: число остовных лесов на $[n]$ из $m$ деревьев, где вершины $1, \dots, m$ обязаны лежать в разных компонентах (корни), равно $m \cdot n^{n - m - 1}$.
\end{itemize}

\textbf{Матричная теорема Кирхгофа (число остовных деревьев произвольного графа).} Пусть дан неориентированный граф (разрешены кратные рёбра; петли сразу выкидываем — в остов они не входят) на $n$ вершинах, где $a_{i,j}$ — число рёбер между $i$ и $j$ (или сумма их весов). Составим \textbf{матрицу Кирхгофа (Лапласа)} $L$ размера $n \times n$: на диагонали стоят степени вершин, вне диагонали — минус кратности рёбер:
\[
  L_{i,i} = \sum_{k \ne i} a_{i,k}, \qquad L_{i,j} = -a_{i,j} \;\; (i \ne j).
\]
Вычеркнем из $L$ \emph{любую} одну строку $r$ и столбец с тем же номером $r$. Определитель оставшейся матрицы $(n - 1) \times (n - 1)$ равен числу остовных деревьев графа (а для взвешенного графа — сумме произведений весов рёбер по всем остовам) и вычисляется Гауссом за $O(n^3)$:
\[
  \tau(G) = \sum_{T \text{ — остов}} \;\prod_{e \in T} w(e) = \det L_{\bar{r}, \bar{r}}.
\]
\emph{Для ориентированного графа (теорема Татта):} чтобы посчитать остовные деревья с корнем $r$, где все рёбра направлены \textbf{от} корня $r$ к листьям, на диагональ ставят \emph{входящие} степени $L^{\mathrm{in}}_{v,v} = \sum_{u \ne v} w_{u \to v}$ (вне диагонали $L^{\mathrm{in}}_{u,v} = -w_{u \to v}$) и вычёркивают строку и столбец корня $r$: ответ $\det L^{\mathrm{in}}_{\bar{r}, \bar{r}}$. Если рёбра направлены \textbf{к} корню $r$, на диагональ ставят \emph{исходящие} степени $L^{\mathrm{out}}_{u,u} = \sum_{v \ne u} w_{u \to v}$.

% ==============================================================================
\section{Уровень 4: Продвинутый арсенал}
% ==============================================================================

\subsection{Инверсия Мёбиуса и биномиальная инверсия}

\textbf{Общая идея инверсий.} Напрямую считать объекты с условием «\textbf{ровно}» ($\gcd = d$, или выполнено ровно $k$ свойств) трудно, зато ослабленное условие «\textbf{делится на $d$}» или «выполнены выбранные $k$ свойств, а на остальные всё равно» обычно считается одной простой формулой или ДП. Инверсии позволяют мгновенно пересчитать массив «ослабленных» ответов в массив «точных».

\textbf{Функция Мёбиуса $\mu(n)$.} Задаёт знаки включений-исключений по простым делителям натурального числа $n = p_1^{\alpha_1} \dots p_r^{\alpha_r}$:
\[
  \mu(n) =
  \begin{cases}
    1, & n = 1, \\
    (-1)^r, & \alpha_1 = \dots = \alpha_r = 1 \text{ ($n$ не делится ни на один квадрат простого)}, \\
    0, & \text{если $n$ делится хотя бы на один квадрат } p_i^2.
  \end{cases}
\]
Предподсчитывается для всех $1 \dots N$ линейным решетом за $O(N)$, мультипликативна ($\mu(ab) = \mu(a)\mu(b)$ при $\gcd(a, b) = 1$) и обладает главным фильтрующим свойством — сумма по делителям обнуляет всё, кроме единицы:
\[
  \sum_{d \mid n} \mu(d) = \iverson{n = 1}, \qquad \sum_{d \mid n} \mu(d)\frac{n}{d} = \varphi(n).
\]

\textbf{Две формы инверсии Мёбиуса (по делителям и по кратным):}
\[
  g(n) = \sum_{d \mid n} f(d) \;\; (\forall n \ge 1) \quad \Longleftrightarrow \quad
  f(n) = \sum_{d \mid n} \mu(d)\, g\!\left(\frac{n}{d}\right) = \sum_{d \mid n} \mu\!\left(\frac{n}{d}\right) g(d);
\]
\[
  G(d) = \sum_{k=1}^{\lfloor M / d \rfloor} F(kd) \;\; (1 \le d \le M) \quad \Longleftrightarrow \quad
  F(d) = \sum_{k=1}^{\lfloor M / d \rfloor} \mu(k)\, G(kd).
\]

\textbf{Как применять это в задачах на массивы и графы с $\gcd$:}
\begin{itemize}
  \item \textbf{Трюк с подстановкой $[\gcd(x, y) = 1] = \sum_{d \mid x,\, d \mid y} \mu(d)$:} условие взаимной простоты двух элементов заменяется суммой по общим делителям $d$, после чего $x$ и $y$ становятся независимыми (достаточно знать для каждого $d$, сколько чисел в массиве кратны $d$):
  \[
    \sum_{i, j} \iverson{\gcd(a_i, a_j) = g} w(i, j) = \sum_{k=1}^{\lfloor M / g \rfloor} \mu(k) \sum_{\substack{i : gk \mid a_i \\ j : gk \mid a_j}} w(i, j).
  \]
  \item \textbf{Переход от «$d \mid \gcd$» к «$\gcd = d$» за $O(M \log M)$:} считаем для каждого $d \in [1, M]$ простую величину $G(d)$ — сколько наборов/путей состоят только из чисел, делящихся на $d$. Тогда точный ответ $F(d)$ для $\gcd = d$ получается либо формулой $F(d) = \sum_{k=1}^{\lfloor M / d \rfloor} \mu(k) G(kd)$, либо даже без функции Мёбиуса обратным проходом от $d = M$ вниз к $1$: $F(d) = G(d) - \sum_{k \ge 2} F(kd)$. Из-за гармонического ряда $\sum_{d=1}^M \lfloor M / d \rfloor = O(M \log M)$ это работает почти линейно.
  \item \textbf{Графы с $\gcd$ весов рёбер:} для каждого $d$ оставляем в графе $H_d$ только рёбра с весом, кратным $d$, и считаем в $H_d$ число путей (через ДП/СНМ) или остовов (Кирхгофом). Каждое ребро попадёт лишь в $\tau(w_e)$ графов ($\tau(x)$ — число делителей, для $x \le 10^6$ не превышает $240$), поэтому суммарный размер всех непустых $H_d$ мал.
\end{itemize}

\textbf{Биномиальная инверсия (переход между «выбрано $k$», «ровно $k$» и «хотя бы $k$»).}
\begin{itemize}
  \item \textbf{Классическая симметричная пара:}
  \[
    b_n = \sum_{k=0}^{n} \binom{n}{k} a_k \quad \Longleftrightarrow \quad a_n = \sum_{k=0}^{n} (-1)^{n - k} \binom{n}{k} b_k.
  \]
  \item \textbf{Рабочая форма для комбинаторного ДП (наддиагональная):} пусть есть $n$ условий $A_1, \dots, A_n$. Важно различать три величины:
  \begin{enumerate}
    \item $E_k$ (\textbf{«ровно $k$»}) — число объектов, где выполнено ровно $k$ условий;
    \item $L_k = \sum_{j=k}^n E_j$ (\textbf{«честное хотя бы $k$»}) — число объектов, где выполнено $\ge k$ условий (каждый посчитан по одному разу);
    \item $G_k = \sum_{|S| = k} \left|\bigcap_{i \in S} A_i\right|$ (\textbf{«зафиксировали $k$ условий»}, в разборах часто нестрого называют «хотя бы $k$») — мы выбрали конкретные $k$ условий, потребовали их, а на остальные $n - k$ не смотрели. Любой объект, у которого на самом деле выполнено $j \ge k$ условий, посчитается в $G_k$ ровно $\binom{j}{k}$ раз!
  \end{enumerate}
  Именно массив $G_0, \dots, G_n$ обычно легко считает ДП. Пересчёт в «ровно $k$» и «честное хотя бы $k$» ($k \ge 1$):
  \[
    G_k = \sum_{j=k}^{n} \binom{j}{k} E_j \quad \Longleftrightarrow \quad
    E_k = \sum_{j=k}^{n} (-1)^{j - k} \binom{j}{k} G_j, \qquad
    L_k = \sum_{j=k}^{n} (-1)^{j - k} \binom{j - 1}{k - 1} G_j.
  \]
  При $k = 0$ формула для $E_0$ совпадает с обычной ФВИ. Все значения $E_0, \dots, E_n$ считаются по $G$ за $O(n^2)$ напрямую или за $O(n \log n)$ одной свёрткой NTT, так как $k! E_k = \sum_{u=0}^{n - k} \frac{(-1)^u}{u!} \cdot (k + u)! G_{k + u}$.
\end{itemize}

% ------------------------------------------------------------------------------
\subsection{Производящие функции (OGF и EGF) + NTT}

\textbf{Как думать про производящие функции.} \textbf{Формальный степенной ряд} — это способ упаковать бесконечную последовательность чисел $a_0, a_1, a_2, \dots$ в один алгебраический объект с переменной $x$. Нам не нужно подставлять вместо $x$ конкретные числа или заботиться о сходимости ряда: переменная $x$ служит просто меткой веса ($x^k \cdot x^m = x^{k + m}$), благодаря чему умножение и функции от рядов автоматически делают за нас комбинаторные свёртки над коэффициентами.

\textbf{Обыкновенные производящие функции (OGF) — для непомеченных объектов.} Последовательности $(a_n)_{n \ge 0}$ (число способов набрать вес $n$ монетами, длина ПСП, число путей) ставится в соответствие ряд $A(x) = \sum_{n=0}^\infty a_n x^n$:
\begin{itemize}
  \item \textbf{Смысл операций:} сложение $A(x) + B(x)$ даёт $a_n + b_n$ (выбор из двух непересекающихся классов); умножение $A(x) B(x)$ даёт свёртку $\sum_{k=0}^n a_k b_{n - k}$ (упорядоченная пара из объекта из $A$ и объекта из $B$ суммарного веса $n$); деление на $(1 - x)$ считает префиксные суммы $[x^n]\frac{A(x)}{1 - x} = \sum_{k=0}^n a_k$; если пустого объекта нет ($a_0 = 0$), то $\frac{1}{1 - A(x)} = 1 + A(x) + A(x)^2 + \dots$ перечисляет произвольные цепочки (последовательности) объектов из $A$.
  \item \textbf{Главные табличные OGF} (обобщённый бином $\binom{\alpha}{n} = \frac{\fall{\alpha}{n}}{n!}$ работает для любого, даже дробного или отрицательного $\alpha$):
  \[
    \frac{1}{1 - x} = \sum_{n=0}^{\infty} x^n, \quad
    (1 + x)^\alpha = \sum_{n=0}^{\infty} \binom{\alpha}{n} x^n, \quad
    \frac{1}{(1 - x)^k} = \sum_{n=0}^{\infty} \binom{n + k - 1}{k - 1} x^n, \quad
    \sum_{n=0}^{\infty} C_n x^n = \frac{1 - \sqrt{1 - 4x}}{2x}.
  \]
\end{itemize}

\textbf{Экспоненциальные производящие функции (EGF) — для помеченных объектов.} Когда элементы пронумерованы метками $1, \dots, n$ (графы, деревья, перестановки), коэффициенты делят на $n!$: $\widehat{A}(x) = \sum_{n=0}^\infty a_n \frac{x^n}{n!}$ (чтобы достать ответ назад, умножаем на $n!$: $a_n = n! [x^n]\widehat{A}(x)$).
\begin{itemize}
  \item \textbf{Почему именно деление на $n!$:} при перемножении $\widehat{C}(x) = \widehat{A}(x)\widehat{B}(x)$ коэффициенты складываются в \emph{биномиальную свёртку} $c_n = \sum_{k=0}^n \binom{n}{k} a_k b_{n - k}$ — это в точности выбор $\binom{n}{k}$ меток под первую фигуру и остальных $n - k$ меток под вторую! Производная сдвигает нумерацию на единицу: $\widehat{A}'(x) = \sum_{n=0}^\infty a_{n + 1} \frac{x^n}{n!}$.
  \item \textbf{Базовые EGF:} $e^x \leftrightarrow a_n = 1$ (одно простое множество на $n$ метках); $\frac{1}{1 - x} \leftrightarrow n!$ (перестановки); $\ln\frac{1}{1 - x} \leftrightarrow (n - 1)!$ ($n \ge 1$, ориентированные циклы); $\frac{e^{-x}}{1 - x} \leftrightarrow !n$ (беспорядки); $\frac{(e^x - 1)^k}{k!} \leftrightarrow \stirlingsecond{n}{k}$ ($k$ непустых кучек).
  \item \textbf{Связь $\exp$ и $\ln$ (переход между связными объектами и всеми объектами):} пусть $\widehat{C}(x) = \sum_{n=1}^\infty c_n \frac{x^n}{n!}$ ($c_0 = 0$) — EGF одного \emph{связного} помеченного объекта (например, связного графа или одного цикла). Тогда $\frac{\widehat{C}(x)^k}{k!}$ — EGF набора из ровно $k$ связных компонент (делим на $k!$, так как порядок самих компонент не важен), а сумма по всем $k \ge 0$ даёт EGF вообще всех объектов $\widehat{G}(x) = \sum_{n=0}^\infty g_n \frac{x^n}{n!}$ ($g_0 = 1$):
  \[
    \widehat{G}(x) = \exp\!\bigl(\widehat{C}(x)\bigr) \quad \Longleftrightarrow \quad \widehat{C}(x) = \ln \widehat{G}(x).
  \]
  \emph{Классический пример:} всего помеченных графов на $n$ вершинах $g_n = 2^{\binom{n}{2}}$ (каждое ребро либо есть, либо нет). Взяв полиномиальный логарифм $\widehat{C}(x) = \ln\!\bigl(\sum_{n=0}^N 2^{\binom{n}{2}} \frac{x^n}{n!}\bigr) \bmod x^{N + 1}$, мы за $O(N \log N)$ сразу получаем количества связных графов $c_1, \dots, c_N$.
\end{itemize}

\textbf{NTT и перемножение многочленов по модулю $998244353$ за $O(N \log N)$.} Чтобы перемножить два многочлена за $O(N \log N)$, их вычисляют в $N = 2^m$ точках (степенях корня из единицы), перемножают поточечно за $O(N)$ и интерполируют обратно. Вместо комплексных чисел и неточного \texttt{double} (FFT) в \textbf{NTT} то же самое делают точно в остатках по простому модулю вида $p = c \cdot 2^K + 1$. Для главного контестного модуля $p = 998244353 = 119 \cdot 2^{23} + 1$ первообразный корень равен $g = 3$ (степени тройки пробегают все остатки $1 \dots p - 1$), поэтому для любой степени двойки $N = 2^m \le 2^{23}$ число $\omega_N = 3^{(p - 1) / N} \bmod p$ является точным корнем степени $N$ из единицы ($\omega_N^N \equiv 1$, $\omega_N^{N/2} \equiv p - 1 \pmod p$):
\[
  \hat{a}_j = \sum_{k=0}^{N - 1} a_k \omega_N^{jk} \bmod p, \qquad
  a_k = N^{-1} \sum_{j=0}^{N - 1} \hat{a}_j \omega_N^{-jk} \bmod p \qquad (0 \le j, k < N).
\]
Оба преобразования считаются бабочкой Кули—Тьюки за $O(N \log N)$ при $N = 2^{\lceil \log_2(n + m - 1) \rceil}$.

\textbf{Операции $\inv, \ln, \exp$ над рядами по модулю $x^N$ за $O(N \log N)$.} Нам нужны только первые $N$ коэффициентов результата ($\bmod\, x^N$, при $N \le p$). Итерация метода Ньютона стартует с $1$ известного коэффициента и на каждом шаге удваивает верную длину $m = 1, 2, 4, \dots$ (суммарное время удвоений $O(1 + 2 + 4 + \dots + N \log N) = O(N \log N)$):
\begin{enumerate}
  \item \textbf{Обратный ряд $B(x) = A(x)^{-1} \bmod x^N$} (существует тогда и только тогда, когда свободный член $a_0 = [x^0]A(x) \not\equiv 0 \pmod p$):
  \[
    B_1(x) = a_0^{-1} \bmod p, \qquad
    B_{2m}(x) \equiv B_m(x)\bigl(2 - A(x)B_m(x)\bigr) \pmod{x^{2m}}.
  \]
  \item \textbf{Логарифм $\ln A(x) \bmod x^N$} (определён при $[x^0]A(x) = 1$, даёт $[x^0]\ln A(x) = 0$; здесь производная и интеграл берутся почленно: $(c_k x^k)' = k c_k x^{k - 1}$, $\int c_k x^k dx = \frac{c_k}{k + 1} x^{k + 1}$):
  \[
    \ln A(x) \equiv \int \Bigl(A'(x) \cdot A(x)^{-1} \bmod x^{N - 1}\Bigr) dx \pmod{x^N}.
  \]
  \item \textbf{Экспонента $\exp A(x) \bmod x^N$} (определена при $[x^0]A(x) = 0$, даёт $[x^0]\exp A(x) = 1$):
  \[
    E_1(x) = 1, \qquad
    E_{2m}(x) \equiv E_m(x)\bigl(1 - \ln E_m(x) + A(x)\bigr) \pmod{x^{2m}}.
  \]
  \item \textbf{Быстрое возведение ряда в любую степень $k$:} выносим младший ненулевой член $A(x) = c x^d \tilde{A}(x)$ (так что $[x^0]\tilde{A}(x) = 1$) и считаем $A(x)^k \equiv c^k x^{kd} \exp(k \ln \tilde{A}(x)) \pmod{x^N}$.
\end{enumerate}

% ------------------------------------------------------------------------------
\subsection{Специальные приёмы}

\textbf{Max-Min Inclusion-Exclusion (переворот максимума в минимумы).} Для любого непустого конечного набора чисел $S$ максимум выражается через минимумы всех непустых подмножеств $T \subseteq S$ по аналогу ФВИ, а по линейности то же самое верно и для матожиданий любых случайных величин $\{X_i\}_{i \in S}$:
\[
  \max(S) = \sum_{\emptyset \ne T \subseteq S} (-1)^{|T| - 1} \min(T), \qquad
  \E\!\left[\max_{i \in S} X_i\right] = \sum_{\emptyset \ne T \subseteq S} (-1)^{|T| - 1} \E\!\left[\min_{i \in T} X_i\right].
\]
\emph{Где это стреляет:} пусть $X_i$ — момент времени, когда впервые произошло $i$-е событие (например, выпал $i$-й купон с неравной вероятностью $p_i$, или закрасился $i$-й бит). Тогда $\max_{i \in S} X_i$ — момент, когда собраны \textbf{все} элементы (считать в лоб тяжело), а $\min_{i \in T} X_i$ — момент, когда сработал \textbf{хотя бы один} элемент из $T$, что является обычным геометрическим ожиданием первого попадания в $T$: $\E[\min_{i \in T} X_i] = \frac{1}{P(\text{за шаг задеть } T)}$.
Обобщение на $k$-й с конца максимум $\max^{(k)}(S)$ ($1 \le k \le |S|$, где $\max^{(1)} = \max$):
\[
  \E\!\left[\max\nolimits^{(k)}_{i \in S} X_i\right] = \sum_{\substack{T \subseteq S \\ |T| \ge k}} (-1)^{|T| - k} \binom{|T| - 1}{k - 1} \E\!\left[\min_{i \in T} X_i\right].
\]

\textbf{Лемма Линдстрёма—Гесселя—Вьенно (LGV — непересекающиеся пути в DAG).} Пусть дан ориентированный граф без циклов (DAG) с весами рёбер $w(e)$ (если просто считаем число наборов путей, ставим $w(e) = 1$; вес одного пути $w(P) = \prod_{e \in P} w(e)$). Обозначим через $e(u, v) = \sum_{P : u \leadsto v} w(P)$ суммарный вес всех путей из $u$ в $v$ (находится обычным ДП по DAG). Выбраны $k$ стартов $A = (a_1, \dots, a_k)$ и $k$ финишей $B = (b_1, \dots, b_k)$. Составим матрицу $M$ размера $k \times k$, где $M_{i,j} = e(a_i, b_j)$. Тогда её определитель считает только системы из $k$ \textbf{не имеющих общих вершин} путей (все пересекающиеся наборы взаимно уничтожаются):
\[
  \det M = \sum_{\sigma \in S_k} \sgn(\sigma) \sum_{\substack{\mathcal{P} = (P_1, \dots, P_k) \\ P_i : a_i \leadsto b_{\sigma(i)},\; V(P_i) \cap V(P_j) = \emptyset}} \prod_{i=1}^{k} w(P_i), \qquad \sgn(\sigma) = (-1)^{\inv(\sigma)}.
\]
\emph{Главный практический случай:} если граф устроен так (например, пути вправо-вверх на клетчатой доске, где старты и финиши идут по порядку вдоль границ), что соединить $a_i$ с $b_{\sigma(i)}$ без пересечения путей возможно \textbf{только напрямую} ($\sigma = \mathrm{id}$, каждый $a_i$ обязан прийти в свой $b_i$), то в сумме остаётся ровно одна перестановка со знаком $+1$, и $\det M$ сразу выдаёт число наборов из $k$ вершинно-непересекающихся путей $a_i \leadsto b_i$. Время: $O(k(|V| + |E|) + k^3)$ (а на доске $M_{i,j}$ вообще считается формулой бинома).

\textbf{Вероятностный метод и рандомизация:}
\begin{itemize}
  \item \textbf{Три опоры вероятностного метода:}
  \begin{enumerate}
    \item \emph{Оценка объединения (Union Bound):} вероятность того, что случится хотя бы один «провал» из $B_1, \dots, B_m$, никогда не превышает суммы их вероятностей: $P(\bigcup_{i=1}^m B_i) \le \sum_{i=1}^m P(B_i)$. Если $\sum P(B_i) < 1$, то с положительной вероятностью ни одно плохое событие не наступит — значит, хороший объект гарантированно существует (и случайная генерация быстро его находит).
    \item \emph{Принцип среднего (первый момент):} случайная величина не может быть всегда строго меньше своего среднего — всегда найдётся исход с $X(\omega_1) \ge \E[X]$ и исход с $X(\omega_2) \le \E[X]$.
    \item \emph{Лемма Шварца—Зиппеля:} ненулевой многочлен $P(x_1, \dots, x_n) \not\equiv 0$ суммарной степени $\le d$ над полем (например, по простому модулю $p$) в случайно выбранной точке $r_1, \dots, r_n \in S$ случайно обращается в $0$ с вероятностью не более $\frac{d}{|S|}$. Это позволяет за секунду проверять равенство сложных алгебраических выражений и определителей матриц с переменными просто подстановкой случайных остатков.
  \end{enumerate}

  \item \textbf{Парадокс дней рождения:} если мы $k$ раз независимо выбираем случайное число из диапазона размера $M$, вероятность ни разу не повториться падает квадратично по $k$:
  \[
    P(\text{нет совпадений}) = \prod_{i=0}^{k - 1}\left(1 - \frac{i}{M}\right) \le \exp\!\left(-\frac{k(k - 1)}{2M}\right), \qquad
    P(\text{есть коллизия}) \le \frac{k(k - 1)}{2M}.
  \]
  Первое совпадение с шансом $\ge 50\%$ ждёт нас уже при $k \approx 1{,}177\sqrt{M}$ ($\Theta(\sqrt{M})$). Главное следствие для хешей: если в \texttt{set}/хеш-таблице лежит или попарно сравнивается $k$ хешей, одиночный модуль $M \approx 10^9$ сломается уже на $k \approx 4 \cdot 10^4$; для $k = 10^6$ обязательно нужен размер пространства $M \gg k^2$ — $64$-битный хеш ($M \sim 2^{64}$) или пара независимых модулей.

  \item \textbf{Zobrist-хеширование (хеш множества без учёта порядка):} каждому виду элемента $x$ (числу массива или паре «клетка—фигура» на доске) один раз сопоставляем случайное $64$-битное число $h(x)$ из \texttt{mt19937\_64} со случайным \texttt{seed}:
  \begin{itemize}
    \item \emph{Для множеств и проверки чётности вхождений:} берём XOR-хеш $H_\oplus(S) = \bigoplus_{x \in S} h(x)$. Он не зависит от порядка элементов, за $O(1)$ пересчитывается при добавлении и удалении, а на подотрезке равен $0$ тогда и только тогда (с ошибкой $2^{-64}$), когда каждое число встретилось чётное число раз.
    \item \emph{Для мультимножеств (важны количества вхождений):} берём обычную сумму $H_+(M) = \sum_{x \in M} h(x) \bmod 2^{64}$ (чтобы проверить совпадение двух мультимножеств или является ли отрезок перестановкой). Для условия «каждое число входит кратно $K$ раз» поддерживаем остатки количеств по модулю $K$ с весами $h(x)$ (или складываем случайные векторы покомпонентно по модулю $K$).
  \end{itemize}

  \item \textbf{Алгоритм Фрейвалдса (проверка равенства матриц $A \cdot B = C$ за $O(k n^2)$):} вместо честного умножения матриц $n \times n$ за $O(n^3)$ генерируем случайный столбец из нулей и единиц $r \in \{0, 1\}^n$ и проверяем равенство векторов, умножая справа налево:
  \[
    A \cdot (B r) \stackrel{?}{=} C r \qquad \bigl(\text{три умножения матрицы на вектор, каждое за } O(n^2)\bigr).
  \]
  Если $AB = C$, равенство выполнится всегда. Если же $AB \ne C$, в матрице $AB - C$ есть хотя бы одна ненулевая строка, и случайный вектор $r$ случайно даст ноль с вероятностью $\le 1/2$. Повторив проверку $k = 35\text{--}40$ раз с независимыми векторами $r$, мы уменьшаем шанс не заметить ошибку до $\le 2^{-k} < 10^{-11}$ за время $O(k n^2)$.
\end{itemize}

\end{document}
